% Photo -> print-ready: page 1 = input/photo-1.jpeg, page 2 = input/photo-2.jpeg (diagram redrawn in TikZ)
% Every word transcribed as printed; doubts are listed in verify.md (NOT silently "fixed").
\documentclass[11pt]{article}
\usepackage[tiro]{edtech}
\geometry{left=20mm,right=24mm,top=14mm,bottom=16mm}
\pagestyle{empty}
\cqstyle{board}
\begin{document}
%% ---------------------------------------------------------------- photo-1
\boardcqheader{name=মানুষ, center=১৮৫, serial=১২৪৬/২০২৫, set=০৪ সেট,
  subject=উৎপাদন ব্যবস্থাপনা ও বিপণন, paper=প্রথম পত্র, syllabus={[২০২৫ সালের সিলেবাস অনুযায়ী]},
  time={২ ঘণ্টা ৩০ মিনিট}, marks=৭০, code=286,
  note={যে কোনো সাতটি প্রশ্নের উত্তর দিতে হবে।}}

\cq{জনাব বাবুল যশোরের গদখালীতে ক্রয়কৃত নিজস্ব ৫ একর জমিতে বিভিন্ন ধরনের ফুল চাষ করেন। তিনি উৎপাদিত ফুল বিভিন্ন দিবসে ব্যবহারের জন্য বিভিন্ন মাধ্যমে ঢাকাসহ দেশের বিভিন্ন শহরে বিক্রির ব্যবস্থা করে লাভবান হন। চলতি বছরে ফুলের যথেষ্ট চাহিদা ও বাম্পার ফলন হলেও নিজস্ব পরিবহণ ব্যবস্থা না থাকায় তিনি যথা সময়ে ফুল বিক্রয় করতে ব্যর্থ হন। ফলে যথেষ্ট লোকসানের সম্মুখীন হন। বিষয়টি নিয়ে তিনি বেশ চিন্তিত।}
{উৎপাদন কী?}
{আর্থিক উৎপাদনশীলতা কীভাবে নির্ণয় করা যায়?}
{জনাব বাবুলের কার্যক্রম উৎপাদনের কোন খাতের অন্তর্ভুক্ত? ব্যাখ্যা কর।}
{উপযোগের ধরণ বিবেচনায় জনাব বাবুলের সমস্যা সমাধানে তোমার সুপারিশ তুলে ধর।}

\cq{২০২৪ সালে মীম গার্মেন্টসে এর উৎপাদিত পোশাকের মূল্য ছিল ১০,০০,০০,০০০ টাকা। উক্ত পোশাক উৎপাদনের জন্য প্রতিষ্ঠানটি কাপড় ও কাঁচামাল বাবদ ৪,০০,০০,০০০ টাকা, শ্রম বাবদ ১,০০,০০,০০০ টাকা, যন্ত্রপাতি বাবদ ২,০০,০০,০০০ টাকা এবং অন্যান্য খরচ বাবদ ১,০০,০০,০০০ টাকা ব্যয় করেন। প্রতিষ্ঠানটির ২০২৩ সালের যন্ত্রপাতির উৎপাদনশীলতা ছিল ৬।}
{উৎপাদনশীলতা কী?}
{রূপগত উপযোগ কীভাবে সৃষ্টি হয়? ব্যাখ্যা কর।}
{২০২৪ সালের ‘মীম গার্মেন্টেসের’ মোট উৎপাদনশীলতা নির্ণয় কর।}
{‘মীম গার্মেন্টেসের’ যন্ত্রপাতির উৎপাদনশীলতার প্রেক্ষিতে প্রতিষ্ঠানটির বর্তমান অবস্থার মূল্যায়ন কর।}
\continued

%% ---------------------------------------------------------------- photo-2
\newpage\setcounter{etq}{0}
\boardcqheader{name=যমুনা, center=১৮২, serial=১১২৬/২০২৬, set=০২ সেট,
  subject=রসায়ন, paper=প্রথম পত্র, syllabus={[২০২৬ সালের সিলেবাস অনুযায়ী]},
  time={২ ঘণ্টা ৩৫ মিনিট}, marks=৫০, code=176,
  instr={ডান পাশের সংখ্যা প্রশ্নের পূর্ণমান জ্ঞাপক। প্রদত্ত উদ্দীপকগুলো মনোযোগসহকারে পড় এবং সংশ্লিষ্ট প্রশ্নগুলোর যথাযথ উত্তর দাও।},
  note={যে কোনো পাঁচটি প্রশ্নের উত্তর দিতে হবে।}}

\cq{একাদশের ছাত্র, আহনাফ ল্যাবরেটরিতে টাইট্রেশনের মাধ্যমে প্রমাণ \ce{Na2CO3} দ্রবণের সাহায্যে অক্সালিক এসিড দ্রবণের ঘনমাত্রা নির্ণয় করার জন্য \SI{10}{mL} \ce{Na2CO3} নিলেন। অন্যদিকে ওয়াসিফ অন্য একটি পরীক্ষণে নমুনা হিসেবে ৮০ মিলিগ্রাম পরিমাণ নিয়ে নির্দিষ্ট লক্ষ্য অর্জনের জন্য দহন প্রক্রিয়া সম্পন্ন করেন।}
{দ্রবণ তাপ কাকে বলে?}
{সাধারণ তাপমাত্রায় \ce{CO2} গ্যাস, কিন্তু \ce{SiO2} কঠিন কেন? ব্যাখ্যা করো।}
{উদ্দীপকে উল্লিখিত আহনাফের পরীক্ষণটি সম্পন্ন করতে ব্যবহৃত যন্ত্রপাতির মধ্যে যেকোনো দুটির ব্যবহারের কৌশল বর্ণনা করো।}
{উদ্দীপকে উল্লিখিত দুটি বিশ্লেষণীয় পদ্ধতির মধ্যে কোনটি পরিবেশবান্ধব? যুক্তিসহ বিশ্লেষণ করো।}

\cq{\vspace{-2pt}\begin{center}\begin{tikzpicture}[font=\normalsize,line width=0.9pt]
  \newcommand\beaker[3]{% x, label, contents
    \draw[rounded corners=7pt] (#1,2.3)--(#1,0)--(#1+2.6,0)--(#1+2.6,2.3);
    \node[align=center] at (#1+1.3,1.15){#3};
    \node at (#1+1.3,-0.35){#2};}
  \beaker{0}{পাত্র : ১}{$4.2\times10^{-3}$ M\\ \SI{50}{mL}\\ $\mathrm{X^{3+}}$ এর দ্রবণ}
  \beaker{3.3}{পাত্র : ২}{$1.5\times10^{-4}$ M\\ \SI{30}{mL}\\ YOH এর\\ জলীয় দ্রবণ}
  \beaker{6.6}{পাত্র : ৩}{পাত্র : ১\\ $+$\\ পাত্র : ২ এর\\ মিশ্রিত দ্রবণ}
\end{tikzpicture}\end{center}
\hspace*{3.2cm}\begin{tabular}{@{}l@{}}Y এর পারমাণবিক সংখ্যা 11\\
$\mathrm{X(OH)_3}$ এর দ্রাব্যতা গুণাংক $1\times10^{-33}\,\mathrm{mol^4L^{-4}}$\end{tabular}}
{সাসপেনশন কাকে বলে?}
{বিশুদ্ধ পানির pH মান 7 কেন?}
{উদ্দীপকের ২ নং পাত্রের দ্রবণের ক্যাটায়নটির শনাক্তকারী পরীক্ষা রাসায়নিক সমীকরণসহ বর্ণনা করো।}
{৩নং পাত্রে অধঃক্ষেপ পড়ার সম্ভাবনা গাণিতিকভাবে বিশ্লেষণ করো।}
\continued
\end{document}
